\section{物理仿真与 AI 的融合}

\subsection{三个世界之间的鸿沟}

访谈中有一段关于"三个世界"的精彩讨论：（1）真实的世界；（2）我们肉眼能观测到的世界；（3）我们用物理模型建模出来的世界。

\begin{importantbox}{三个世界之间的 Gap}
\textbf{第一个 Gap}（真实世界 vs 可观测世界）：每个人只能观测到世界极小的一部分，所以每个人的认知都有巨大局限性。这个 gap 的本质是\textbf{观测能力的有限性}。

\textbf{第二个 Gap}（可观测世界 vs 仿真世界）："差了十万八千里"。但原因不同——这个 gap 是因为我们\textbf{没有足够的 compute 和 data} 去建模现实世界。在某些经典场景（如 CFD 与风洞测试的关联性）已经很好了，但"飞机起飞前肯定还要吹风洞，你不能说 CFD 就结束了"。
\end{importantbox}

他对此的感叹充满哲学意味："挺无奈的。我们作为一个物种也许运行在别人的一个模拟器里面，我们居然还妄想造一个自己的模拟器能够模拟我们自己。"但他同时认为，GPT-4 这样的模型能 mimic 人类大脑的运作——"已经是一个奇迹了"。

\subsection{Video Model vs. Physical Simulation}

"我们为什么还要用太极做 rendering、做 physical simulation？为什么不用 Sora？"——这个问题直指 Graphics 学科在 AI 时代的存在价值。

胡渊鸣的回答很务实："五年前99行写冰雪奇缘，现在给 Sora 的 prompt 只用写20个单词。有的场景确实革了 Graphics 人的命。但 Graphics 在做的很多事情促成了 Sora 的诞生。"

最终的演变方向是"\textbf{learning 为主}"，但 simulator 会以某种形式存活——作为 inductive bias、synthetic data 来源，或 learnable module 的一部分。Neural Graphics（DLSS、NeRF 等）已经在逐步替代传统 ray tracing，而太极 2.0 的定位正是"supporting Graphics and AI"。

\subsection{Sim-to-Real Gap：仿真的根本挑战}

胡渊鸣指出，所有现有 simulator（Isaac、MuJoCo、Drake、PyBullet、Genesis）都面临一个共同问题：\textbf{sim-to-real gap}。

\begin{warningbox}{物理定律是准的，但边界条件不是}
牛顿第二定律是准的，相对论也是准的。但很多 boundary condition 是不准的——比如摩擦模型，光滑表面和粗糙表面的摩擦机制完全不同，一块玻璃放在湿漉漉的表面上的摩擦又不一样。这些建模全都是不准的。
\end{warningbox}

他认为未来的 simulator 需要变得更加 data-driven，与神经网络结合。这不意味着要用 Transformer 重新学习牛顿第二定律（本来可以用几个 FLOPS 做完的事情，用几十亿 FLOPS 去做没有意义），而是要让 data-driven 方法去弥补那些\textbf{近似模型无法精确描述}的部分，比如材料的 stress-strain curve。

\subsection{Inductive Bias：物理先验在 AI 中的价值}

胡渊鸣特别强调了 inductive bias（归纳偏置）在 AI 系统中的重要性：

\begin{importantbox}{3D Inductive Bias 的效率}
以 video generation 为例：camera projection 在传统图形学中只需约50个 FLOPS，但如果让 AI 来学习这个映射，可能要加上 billion 级别的参数。所以，camera projection 应该被 hardcode，把模型的容量留给真正需要学习的部分。类似地，video model 中 camera 转360度后看到完全不一样的世界——如果有基本的 3D 表示，至少可以解决 physical persistency 问题。
\end{importantbox}

他总结了一个精辟的平衡原则：

\begin{quote}
"Simulator 里面需要有更多的 AI，AI 里面也需要更多的 3D inductive bias。要找到一个恰到好处的平衡点——不要越俎代庖，不要低估模型的学习能力，但也不要把所有事情交给 learning。"
\end{quote}

\subsection{Synthetic Data 的局限性}

关于 Sergey Levine 在 "The Spork of AGI" 博客中对 simulation 的抵触——认为 simulation 反而是 robotics foundation model 的限制，应该用 real world data——胡渊鸣表示"有道理"：

\begin{warningbox}{Synthetic Data 不能作为主要数据源}
"如果所有 data 都是 synthetic data，或者大部分 data 是 synthetic data，一定不会 work。因为能够 synthesize 的 data 一定 follow 某些 rules，这些 rules 一定是人为思考出来的。现实世界中一亿种情况里，rules 可能只能 cover 一万条。"——但 synthetic data 可以用来 enhance 一个已有模型，而非 build from scratch。
\end{warningbox}

\subsection{模拟的极限与嵌套宇宙}

在讨论"物理仿真的极限"时，胡渊鸣给出了一个令人沉思的回答。如果我们自己的世界就是被模拟出来的，那么\textbf{模拟器里的模拟器不会比母体模拟器更强}。

\begin{knowledgebox}{嵌套模拟的计算复杂度论证}
"也许你在我们的世界里面去造一个模拟器，模拟整个世界并使其运行速度比现实时间更快——这有可能是 unlikely 的。除非你只做 small scale 的模拟。这也意味着\textbf{预测股市可能是很难的}，因为股市涉及整个世界所有人都参与其中——你需要模拟的系统和现实世界一样大。但你做一个风洞模拟——这是可以做的。"
\end{knowledgebox}

这个洞察将小时候"桃树是否在我不看的时候生长"的哲学思考与成年后的计算复杂度理论优雅地连接起来。

\subsection{LLM 的柏拉图洞穴}

被问到"GPT 这样的 LLM 是不是生活在柏拉图所说的洞穴里"时，胡渊鸣的回答出人意料：

\begin{quote}
"是。因为它的 loss function 就是它的 training data，所做的一切就是 minimize 这个 loss function。但——我们每个人难道不是生活在柏拉图的洞穴里面吗？我们每个人看到的东西可能比 LLM 还要少。只不过 LLM 有时候犯常识性错误，所以你经常笑它。但我们真的就比 LLM 强吗？我觉得不一定。"
\end{quote}

他认为 AGI 可以从现有的 text、video、image data 中诞生，不一定需要人类对物理世界的全部认知"内建"进去。关键是给 AGI \textbf{tool use} 能力——让它拥有自己的计算机，能写代码去分析物理现象。"它可以自己发现自己的牛顿定律。"

\subsection{如果重新做一个 Simulator}

如果今天重新做一个 simulator，胡渊鸣会以 \textbf{learning 为主驱动}，而不是牛顿定律为主。他会保留的"底层元素"包括：
\begin{itemize}
\item 物体不会凭空消失（physical persistency）
\item 两个物体不能占据同一个位置（collision）
\item Camera 的 perspective projection
\end{itemize}

而将所有 physical law（包括弹塑性形变、流体动力学等）全部搞成 data-driven。最终得到的系统既不是传统 simulator，也不是纯 neural network，而是\textbf{"嵌入了 neural network 的 3D inductive bias"}——一种让 neural network 更容易学习的结构化先验。

\subsection{AGI 需要物理学吗}

胡渊鸣认为 AGI "可能不特别需要牛顿定律和拉格朗日方程"，但需要 \textbf{tool use}——AGI 需要有自己的计算机，能写 finite element 代码去分析物理现象。

\begin{knowledgebox}{Next Token Prediction 的本质局限}
不管 token 信息量多大多小，Transformer 的每个 token 都要花费固定的 compute（$2 \times$ parameter count / sparsity via MoE）。但有些 token 天然需要比其他 token 更多的 compute——这些额外 compute 不能发生在 LLM 当前范式内部，而是要让 LLM 去操作外部工具。所以 AGI 不需要人类对物理世界的全部认知"内建"进去，它可以"自己发现自己的牛顿定律"。
\end{knowledgebox}

这个观点有一个重要推论：当前 scaling up Transformer 的方式（增加参数量，让每个 token 的 compute 更多）本质上是在用暴力解决一个架构层面的问题。更聪明的做法可能是——像太极那样——让不同粒度的计算用不同的方式处理。

\subsection{关于 Robotics 的思考}

尽管胡渊鸣的技术背景与 Robotics 高度相关（物理仿真、可微编程、3D 生成），他坦言对 Robotics"兴趣程度没有那么高"。

\begin{quote}
"我觉得 Robotics 是一个 impact 非常大的事情。可能唯一一个我没有开始做 Robotics 的原因是——我自己今天对它的兴趣程度没有那么高。我相信我只有做自己喜欢的事情才能做好。"
\end{quote}

但他也承认 Robotics 与 Meshy 有潜在的交叉点：很多合作伙伴已经在用 Meshy 生成茶杯、碗、桌子等3D资产用于 Robotics 训练。未来如果 Meshy 支持 articulated objects（有关节的物体），这个交叉将更加显著。

他对 Robotics 中 simulation 角色的判断也很有洞见：让 robot 操作桌上的蔬菜和盘子，\textbf{比飞机上天更难}——"因为天上无非就是空气，但当你要更细致地 manipulate 桌上的东西，各种边界条件的不一致会导致 simulation 在现实中非常难落地。"

\subsection{本章小结}

物理仿真正处于从 rule-based 向 data-driven 转型的关键时期。传统仿真的瓶颈不在物理定律本身，而在边界条件的精确建模。未来的方向是保留底层的 3D 先验（作为 inductive bias），让 learning 处理复杂的、难以精确建模的部分。太极 2.0 正朝这个方向演进。

